\section{Heroes and Archetypes}
\label{sec:heroes-system}

\subsection{Overview}

Heroes are the core units the player builds, equips and sends into Maps. 
Each Hero belongs to an \textbf{Archetype} which defines their role, core stats, AI behaviour and passive traits. 
Heroes do not have traditional skill trees; instead, they use the Skill System via sockets, modules and orbits 
(see Section~\ref{sec:skill-system}).

\begin{itemize}
	\item \textbf{Archetype} = role, stat focus, base behaviour, passive identity.
	\item \textbf{Hero} = instance of an archetype with a level, gear, skill modules, orbits, runes and crowns.
	\item \textbf{Team} = up to 4 Heroes fielded together in combat.
\end{itemize}

\subsection{Team Size and Hero Unlocks}

\subsubsection*{Team Size}

\begin{itemize}
	\item The player starts with a \textbf{single Hero slot} (team size 1).
	\item Additional team slots are unlocked over time (via progression or research) up to a maximum of \textbf{4 Heroes per team}.
	\item The game is balanced around 1--4 Heroes auto-battling together.
\end{itemize}

\subsubsection*{Hero Unlocks}

Heroes are not all available from the start:

\begin{itemize}
	\item The player starts with exactly \textbf{one archetype} (starter Hero).
	\item Additional Heroes are unlocked in \textbf{packs of 2--3 archetypes}, themed together (e.g. „ranged pack“, „support pack“).
	\item Unlocking and owning are separate:
	\begin{itemize}
		\item \emph{Unlock} makes an archetype available to be recruited.
		\item \emph{Recruitment} may cost a rare currency.
	\end{itemize}
	\item The first few (approx.\ 5) Heroes are granted through story/progression without heavy cost.
	\item Further Heroes are primarily unlocked via the Heroes Research Tree and then purchased with rare currency.
\end{itemize}

\subsection{Archetype Roster}

Table~\ref{tab:archetype-roster} lists the planned archetypes, their primary role and core stat focus. 
This roster is larger than the MVP set; the MVP subset is marked explicitly.

\begin{table}[h]
	\centering
	\begin{tabular}{p{2.8cm}p{3.2cm}p{2.5cm}p{5cm}p{1.5cm}}
		\textbf{Archetype} & \textbf{Role} & \textbf{Core Stat} & \textbf{Identity (Short)} & \textbf{MVP} \\
		\hline
		Vanguard & Tank / Frontline & CON & High durability, control, damage mitigation, „front wall“. & Yes \\
		Piercer & STR DPS / Execute & STR & Physical single-target burst, armor penetration, executes. & Future \\
		Arcanist & Caster / Control & INT & Elemental damage, DoTs, debuffs and control. & Yes \\
		Ranger & Ranged Hybrid & DEX & Mobile ranged attacker, projectiles, kiting tools. & Yes \\
		Marksman & Ranged Single-Target & DEX & Long-range precision, crit-focused, positional bonuses. & Future \\
		Alchemist & DoT / Zone Control & INT & Damage-over-time, ground effects, debuffs. & Future \\
		Summoner & Minion Master & WIL & Summoned allies, scaling with minion count and survival. & Yes \\
		Hierophant & Support / Nature & WIL & Healing, auras, protection, buff duration. & Future \\
		Bruiser & Sustain Melee & STR & Melee DPS with self-sustain and high HP. & Future \\
		Saboteur & Traps / Procs & INT & Traps, delayed damage, proc-based synergies. & Future \\
		Oracle & Buff / Debuff Support & WIL & Team-wide buffs and enemy debuffs, predictive flavour. & Future \\
	\end{tabular}
	\caption{Archetype roster and core roles}
	\label{tab:archetype-roster}
\end{table}

For MVP, the core playable set is: \textbf{Vanguard}, \textbf{Arcanist}, \textbf{Ranger}, \textbf{Summoner}. 
Additional archetypes are planned for later iterations once core systems are stable.

\subsection{Stats and Growth}

All archetypes share the same five core stats:

\begin{itemize}
	\item STR, DEX, CON, INT, WIL (see Section~\ref{sec:skill-system}).
\end{itemize}

Each archetype defines:

\begin{itemize}
	\item 1 primary \textbf{Core stat} (e.g. CON for Vanguard, INT for Arcanist),
	\item 1--2 secondary stats,
	\item remaining stats as low-impact for that archetype.
\end{itemize}

Exact numeric growth per level and per-stat distributions are defined in the balance sheets. 
For the GDD, it is sufficient that:

\begin{itemize}
	\item Vanguard and Bruiser lean heavily into \textbf{CON/STR},
	\item Ranger and Marksman lean into \textbf{DEX} (with some INT or STR depending on design),
	\item Arcanist, Alchemist, Saboteur rely on \textbf{INT},
	\item Summoner, Hierophant, Oracle rely on \textbf{WIL}.
\end{itemize}

\subsection{Archetype Passives}

Each archetype has a small set of defining passive traits. 
These are always active once the Hero is in the team and help differentiate playstyle and build direction.

\subsubsection*{Vanguard (Tank, CON Core)}

\begin{itemize}
	\item \textbf{Bulwark:} +20\% Armor.
	\item \textbf{Stand Your Ground:} reduced impact from slows / stuns and/or can move through enemies without being body-blocked (exact implementation to be decided).
	\item \textbf{First Hit Guard:} the first hit taken from each enemy is reduced by 50\% (subject to implementation feasibility).
\end{itemize}

\subsubsection*{Piercer (STR Core, Execute)}

\begin{itemize}
	\item \textbf{Deadly Precision:} +50\% critical damage multiplier.
	\item \textbf{Hemorrhage:} unlocks a stackable Bleed DoT effect on hits.
	\item \textbf{Merciless:} enemies below 40\% HP take +15\% damage from Piercer.
\end{itemize}

\subsubsection*{Arcanist (INT Core, DoT/Control)}

\begin{itemize}
	\item \textbf{Arcane Echo:} +25\% debuff duration.
	\item \textbf{Overchannel:} +20\% spell damage.
	\item \textbf{Flow of Mana:} skills have -15\% cooldown.
\end{itemize}

\subsubsection*{Ranger (DEX Core, Projectiles/Kiting)}

\begin{itemize}
	\item \textbf{Volley:} +1 projectile count for applicable skills.
	\item \textbf{Through the Line:} +25\% pierce chance (projectiles can hit multiple enemies).
	\item \textbf{Hunter’s Focus:} +15\% range.
\end{itemize}

\subsubsection*{Marksman (DEX Core, Precision)}

\begin{itemize}
	\item \textbf{Headshot:} +20\% critical chance.
	\item \textbf{Longshot:} +20\% damage against targets further than a threshold distance (e.g. 8m).
	\item \textbf{Sniper Discipline:} +10\% attack speed while stationary (after not moving for a short duration, e.g. 5s).
\end{itemize}

\subsubsection*{Alchemist (INT Core, DoTs/Zones)}

\begin{itemize}
	\item \textbf{Virulent Concoction:} DoTs can stack +10 additional times.
	\item \textbf{Toxic Spread:} on kill, existing DoTs can spread to a nearby target.
	\item \textbf{Lingering Poison:} +25\% DoT duration.
\end{itemize}

\subsubsection*{Summoner (WIL Core, Minions)}

\begin{itemize}
	\item \textbf{Broodmaster:} +2 summon count (base summon count = 3).
	\item \textbf{Empowered Minions:} minions have +20\% HP and +15\% damage.
	\item \textbf{Hive Mind:} +5\% damage per living minion.
\end{itemize}

\subsubsection*{Hierophant (WIL Core, Support/Nature)}

\begin{itemize}
	\item \textbf{Blessing Aura:} allies in a small radius (e.g. 4m) regenerate 1\% HP/sec.
	\item \textbf{Sanctuary:} the next attack against nearby allies is negated (cooldown, e.g. every 10s).
	\item \textbf{Spirit Bond:} buffs applied by Hierophant last +25\% longer.
\end{itemize}

\subsubsection*{Bruiser (STR Core, Sustain Melee)}

\begin{itemize}
	\item \textbf{Brawler’s Might:} +20\% attack damage.
	\item \textbf{Blood for Strength:} +10\% life leech factor.
	\item \textbf{Thick Skin:} +15\% maximum HP.
\end{itemize}

\subsubsection*{Saboteur (INT Core, Traps/Procs)}

\begin{itemize}
	\item \textbf{Tinker’s Reserve:} traps have +1 charge.
	\item \textbf{Suspicious Gear:} increased chance to apply additional DoT stacks.
	\item \textbf{Expertise:} increased area of effect for trap-triggered effects.
\end{itemize}

\subsubsection*{Oracle (WIL Core, Buff/Debuff)}

\begin{itemize}
	\item \textbf{Prophetic Weakness:} enemy debuffs last +25\% longer.
	\item \textbf{Battle Clairvoyance:} allies in a radius (e.g. 6m) gain +10\% critical chance.
	\item \textbf{Foretold Resilience:} allies in a radius gain +10\% elemental resistance.
\end{itemize}

% NOTE: Numerical values are initial design targets and will be tuned in the balance phase.

\subsection{Sockets and Orbit Progression (Hero Perspective)}

Details of sockets, modules and orbits are defined in the Skill System (Section~\ref{sec:skill-system}). 
From the Heroes perspective, the important points are:

\begin{itemize}
	\item Each Hero has a fixed set of \textbf{4--5 skill sockets} (chakra-style, not tied to gear).
	\item The Hero starts with \textbf{1 active socket} at level 1.
	\item Additional sockets unlock as the Hero levels up:
	\begin{itemize}
		\item all sockets should be unlocked by approximately \textbf{level 60},
		\item the remaining levels (60--100) primarily increase orbit progression and prepare for the Crown system.
	\end{itemize}
	\item Each level grants orbit points that can be spent on socket orbits; exact points-per-level and bonus sources are defined in the Skill/Balance sections, but the intent is:
	\begin{itemize}
		\item base: around 1--2 orbit points per level,
		\item additional points from Map milestones and socket set completion,
		\item a „pure“ level 100 build should be able to complete the intended orbit path.
	\end{itemize}
\end{itemize}

Level 100 marks the end of regular Hero progression: after that, only Crowns and meta systems extend power.

\subsection{AI Profiles (Behaviour)}

Heroes are fully AI-driven in combat. The player does not manually set skill priorities. 
Instead, each archetype has a characteristic AI profile:

\begin{itemize}
	\item \textbf{Vanguard / Bruiser:} 
	\begin{itemize}
		\item prioritise enemies that threaten allies or are closest to the front,
		\item use defensive and control skills early in fights,
		\item stay near the main enemy cluster.
	\end{itemize}
	\item \textbf{Piercer / Marksman / Ranger:}
	\begin{itemize}
		\item prioritise low-HP or high-value targets,
		\item maintain distance when possible (Ranger/Marksman),
		\item use burst skills on high-HP targets or bosses.
	\end{itemize}
	\item \textbf{Arcanist / Alchemist / Saboteur:}
	\begin{itemize}
		\item prioritise clustering (AoE) and debuff application,
		\item use control skills on elites/bosses when available,
		\item Saboteur places traps ahead of enemy movement paths.
	\end{itemize}
	\item \textbf{Summoner:}
	\begin{itemize}
		\item maintains minion count as a priority,
		\item uses skills that buff or reposition minions when possible,
		\item stays behind frontline Heroes.
	\end{itemize}
	\item \textbf{Hierophant / Oracle:}
	\begin{itemize}
		\item prioritise healing and shielding low-HP allies,
		\item apply buffs early in encounters and keep them up,
		\item apply key debuffs to high-HP or elite targets.
	\end{itemize}
\end{itemize}

Exact AI rules, thresholds and target selection weights are part of the Technical Design Document, 
but this section defines the intended behaviour per archetype from a design perspective.

\subsection{MVP Hero Set}

For the MVP, the target is:

\begin{itemize}
	\item \textbf{4 core archetypes} fully implemented:
	\begin{itemize}
		\item Vanguard (tank),
		\item Arcanist (full caster),
		\item Ranger (bow/ranged fighter),
		\item Summoner (minion-focused).
	\end{itemize}
	\item Each with:
	\begin{itemize}
		\item a functional set of 2--3 key Skill Modules,
		\item full socket unlocks by around level 60,
		\item a meaningful orbit layout (Tier I/II, with Tier III and Crowns optional as late/endgame).
	\end{itemize}
	\item Additional archetypes (Marksman, Alchemist, Hierophant, Bruiser, Saboteur, Oracle) 
	are planned as \textbf{post-MVP extensions}.
\end{itemize}

% TODO: Link to concrete archetype balance sheets (stat spreads, module lists, orbit designs) once available.
